\documentclass[11pt]{article}

\usepackage[a4paper,margin=25mm]{geometry}
\usepackage{amsmath,amssymb,amsthm}
\usepackage{booktabs}
\usepackage{microtype}
\usepackage[hidelinks]{hyperref}

\newtheorem{proposition}{Proposition}
\newcommand{\ReLU}[1]{\left[#1\right]_{+}}
\newcommand{\SARwb}{\mathrm{SAR}_{\mathrm{wb}}}
\newcommand{\APDmax}{\mathrm{APD}_{\mathrm{max}}}

\title{Yaw in deterministic human exposure models\\
\large An engineering note}
\author{Working note}
\date{\today}

\begin{document}
\maketitle

\begin{abstract}
Absolute compass direction should not affect human RF exposure when the source
and propagation model has no preferred horizontal direction. This does not mean
that rotating a body inside one fixed street or one fixed set of multipath
components leaves exposure unchanged. The useful engineering result is more
specific. An axisymmetric local field gives exact yaw independence. An
isotropic city model gives yaw independence across the population or ensemble.
When body heading is unknown and uniform, its effect on whole-body SAR can be
averaged exactly without rerunning propagation. This note states these claims,
their limits, and what they imply for deterministic exposure studies.
\end{abstract}

\section{The claim in engineering terms}

There are three different statements that are easy to conflate.

\begin{enumerate}
  \item \textbf{Compass neutrality.} Rotating the complete scene, the body, and
  every field vector together cannot change the exposure result. North is only
  a coordinate label.

  \item \textbf{Yaw independence in one scene.} Rotating the body while keeping
  one street, base station, or dominant path fixed can change exposure. This is
  a physical change in relative orientation.

  \item \textbf{Population yaw independence.} If body heading is independent of
  the local field and has no preferred horizontal direction, population
  exposure does not depend on the compass yaw chosen by the modeler.
\end{enumerate}

The first statement is always required of a physically consistent model. The
second is generally false in a deterministic city realization. The third is
the relevant statement for a compass-neutral population.

No symmetry of the human body is needed. The body may be anatomically
asymmetric and may have different tissue distributions on its front, back, and
sides. The symmetry assumption belongs to the environment or to the heading
distribution.

\section{A compact deterministic statement}

Represent the local incoherent field by multipath components

\begin{equation}
  \mathcal{P}=\{(S_i,\widehat{\mathbf{k}}_i)\}_{i=1}^{N},
\end{equation}

where $S_i$ is incident power density and
$\widehat{\mathbf{k}}_i$ is the propagation direction. Let
$D(B,\mathcal{P})$ be a scalar endpoint computed for body $B$. Examples are
absorbed power and whole-body SAR. A physically consistent endpoint obeys

\begin{equation}\label{eq:common-rotation}
  D(R_\psi B,R_\psi\mathcal{P})=D(B,\mathcal{P}),
\end{equation}

where $R_\psi$ is a yaw rotation about the vertical axis. This equation only
says that rotating the coordinate system changes nothing.

\begin{proposition}[Exact yaw independence]
If the local incident field is unchanged by every yaw rotation,

\begin{equation}\label{eq:axisymmetric-field}
  R_\psi\mathcal{P}=\mathcal{P}
  \quad\text{for all }\psi,
\end{equation}

then $D(R_\psi B,\mathcal{P})=D(B,\mathcal{P})$ for every body yaw.
\end{proposition}

\begin{proof}
Use Eq.~\eqref{eq:axisymmetric-field} to rotate the field together with the
body, then apply Eq.~\eqref{eq:common-rotation}.
\end{proof}

Equation~\eqref{eq:axisymmetric-field} does not require equal power at every
elevation. Base stations may be above the person, ground reflections may arrive
from below, and antenna downtilt may be retained. Only horizontal azimuth is
removed.

This exact proposition is deliberately strong. A finite set of paths from one
ray-traced site will almost never satisfy it. It is best understood as the
limit reached after azimuthal averaging, not as a description of every street.

\section{What flat ground does and does not prove}

An infinite flat homogeneous ground plane has no preferred horizontal
direction. Its reflection coefficient may depend on frequency, material,
polarization, and elevation angle without introducing a compass direction.
Flat ground therefore handles the ground part of the symmetry argument.

It does not make a complete deterministic scene yaw-independent. A wall, street
canyon, antenna sector, or single line-of-sight path can still select one
horizontal direction. Flat ground is sufficient only when the remaining source
and propagation description is also horizontally symmetric.

A weaker and more realistic assumption is statistical isotropy. Suppose base
station locations, heights, powers, antenna marks, and multipath directions
have the same joint distribution after any common yaw rotation. Then

\begin{equation}\label{eq:distributional-result}
  D(R_\psi B,\mathcal{P})
  \overset{d}{=}
  D(B,\mathcal{P}).
\end{equation}

This is an ensemble statement. A single realization can remain strongly
yaw-sensitive. For example, one dominant line-of-sight component has a uniform
azimuth across many sites, but it still illuminates a particular side of the
body at any one site.

The minimal reasonable assumptions for Eq.~\eqref{eq:distributional-result}
are:

\begin{itemize}
  \item the exposure solver is covariant under a common rotation

  \item the joint source and propagation law has no preferred horizontal
  direction

  \item any correlation between body heading and the environment is either
  absent or modeled explicitly
\end{itemize}

Poisson base stations, independent paths, equal heights, and a symmetric human
phantom are not required.

\section{Unknown yaw can be removed exactly for whole-body SAR}

For a fixed deterministic field, let the body yaw
$\Theta\sim\mathrm{Uniform}(0,2\pi)$. Whole-body absorbed power and whole-body
SAR are linear in incoherent path power. Therefore the yaw mean can be computed
either by averaging body orientations or by spreading every path uniformly in
azimuth:

\begin{equation}\label{eq:yaw-average}
  \mathbb{E}_{\Theta}
  \left[\SARwb(R_\Theta B,\mathcal{P})\right]
  =
  \SARwb(B,\overline{\mathcal{P}}),
\end{equation}

where $\overline{\mathcal{P}}$ is the azimuthally averaged angular power
distribution. Propagation is run once. Only the body coupling is averaged.

At AEGIS level 2 this average has a closed form for every triangle and path.
Let $\widehat{\mathbf{n}}$ be a body-surface normal and let
$\mathbf{q}=-\widehat{\mathbf{k}}$ be the direction used in the incidence law.
Define

\begin{equation}
  a=n_zq_z,
  \qquad
  b=\lVert\mathbf{n}_{\mathrm{h}}\rVert
    \lVert\mathbf{q}_{\mathrm{h}}\rVert.
\end{equation}

The exact uniform-yaw mean of the incidence factor is

\begin{equation}\label{eq:closed-form}
  F(a,b)=
  \frac{1}{2\pi}\int_0^{2\pi}\ReLU{a+b\cos\psi}\,\mathrm{d}\psi
  =
  \begin{cases}
    \ReLU{a}, & b=0,\\
    0, & a\leq-b,\\
    a, & a\geq b,\\
    \dfrac{a\arccos(-a/b)+\sqrt{b^2-a^2}}{\pi}, & |a|<b.
  \end{cases}
\end{equation}

This is the main computational result. After yaw averaging, the coupling
depends on path elevation and surface slope, but not on either azimuth. A
dedicated semantic-twin coupling endpoint implements
Eq.~\eqref{eq:closed-form} exactly at AEGIS level 2. The current production
campaign does not call that endpoint. It uses fixed route-tangent yaw.

\subsection{The important exception for peak APD}

The same shortcut cannot be applied blindly to $\APDmax$. In general,

\begin{equation}\label{eq:peak-warning}
  \max_{\mathbf{r}}
  \mathbb{E}_{\Theta}[S_{\mathrm{ab}}(\mathbf{r},\Theta)]
  \neq
  \mathbb{E}_{\Theta}
  \left[\max_{\mathbf{r}}S_{\mathrm{ab}}(\mathbf{r},\Theta)\right].
\end{equation}

The first quantity is the peak of the yaw-mean surface field. The second is the
mean of the peak exposure experienced at each yaw. The second quantity is the
natural population endpoint if each person has one unknown yaw. It requires
evaluating the surface maximum for each sampled orientation, although the
propagation paths still need to be computed only once.

\section{Consequences for deterministic modeling}

The correct treatment of yaw depends on the estimand.

\begin{table}[ht]
\centering
\caption{Recommended yaw treatment by study target.}
\label{tab:recommendations}
\begin{tabular}{p{0.34\linewidth}p{0.58\linewidth}}
\toprule
Study target & Recommended treatment \\
\midrule
One specified route traversal
& Retain route-tangent yaw. Report that the result describes a person facing
the walking direction. \\

Compass-neutral population
& Use the exact uniform-yaw mean for absorbed power and whole-body SAR. For
peak APD, average the per-yaw peaks. \\

One site with unknown heading
& Report the yaw mean and a sensitivity interval, such as the minimum and
maximum or selected quantiles over yaw. \\

Statistically isotropic city ensemble
& A fixed yaw has the correct marginal distribution across ideal independent
realizations. Verify that the simulated site sample is large and balanced
enough for this approximation. \\

Body-fixed phone source
& Absolute compass yaw is irrelevant because the phone rotates with the body.
Phone position and orientation relative to the anatomy remain essential. \\
\bottomrule
\end{tabular}
\end{table}

For the present semantic twin, fixed route-tangent yaw is a legitimate answer
to a narrow question about a prescribed pedestrian traversal. It should not be
presented as a generic population estimate. A stronger result set would report:

\begin{enumerate}
  \item the route-tangent value

  \item the exact uniform-yaw mean for whole-body SAR

  \item a yaw sensitivity range for whole-body SAR and peak APD
\end{enumerate}

This separates information supplied by the scene from an arbitrary heading
choice. If the range is small, detailed azimuth reconstruction adds little to
the endpoint. If the range is large, the field contains a dominant direction
and fixed yaw needs a behavioral justification.

\section{A practical validation experiment}

One numerical experiment is enough to test the relevance of yaw in the current
pipeline.

\begin{enumerate}
  \item Select representative standpoints with line-of-sight, obstructed, and
  multipath-rich conditions.

  \item Keep the traced paths fixed and rotate the body through 24 yaw angles.

  \item For whole-body SAR, compare the numerical mean with the exact result in
  Eq.~\eqref{eq:closed-form}.

  \item Report the mean, coefficient of variation, minimum, and maximum. For
  peak APD, also report the mean of the per-yaw peaks.
\end{enumerate}

A short harmonic summary can be useful as a diagnostic, but it need not become
the paper's mathematical center. The zero-order azimuthal component gives the
yaw mean. The remaining components measure yaw sensitivity. In practice, the
24-angle sweep communicates the same result more directly.

\section{Limits}

The argument above assumes incoherent power addition. For coherent fields, the
full complex field and its spatial coherence must transform consistently.
Azimuthal isotropy of scalar path powers alone is not enough.

Polarization can be retained, but its joint directional distribution must also
be horizontally rotation invariant for the ensemble claim. Sloped terrain,
fixed antenna sectors, aligned street grids, and correlations between walking
direction and street direction are genuine sources of anisotropy. They should
be modeled as such.

The central engineering conclusion is modest. Absolute north is not an exposure
variable. Relative body-to-field orientation can be. When that relative yaw is
unknown and unstructured, whole-body SAR should be marginalized over it rather
than fixed by convenience.

\end{document}
